\documentclass{article}

\PassOptionsToPackage{numbers, sort, compress}{natbib}
\usepackage[main,final]{neurips_2025}

\usepackage[T2A]{fontenc}
\usepackage[utf8]{inputenc}
\usepackage{hyperref}
\usepackage{url}
\usepackage{booktabs}
\usepackage{nicefrac}
\usepackage{microtype}
\usepackage{xcolor}

%%%

\usepackage{subcaption}
\usepackage{graphicx}
\usepackage{multirow}
\usepackage{amsmath,amssymb,amsfonts}
\usepackage{amsthm}
\usepackage{mathrsfs}
\usepackage{xcolor}
\usepackage{textcomp}
\usepackage{manyfoot}
\usepackage{algorithm}
\usepackage{algorithmicx}
\usepackage{algpseudocode}
\usepackage{listings}

\newtheorem{theorem}{Theorem} % continuous numbers
%%\newtheorem{theorem}{Theorem}[section] % sectionwise numbers
%% optional argument [theorem] produces theorem numbering sequence instead of independent numbers for Proposition
\newtheorem{proposition}[theorem]{Proposition}% 
\newtheorem{lemma}{Lemma}% 
%%\newtheorem{proposition}{Proposition} % to get separate numbers for theorem and proposition etc.

\newtheorem{example}{Example}
\newtheorem{remark}{Remark}

\newtheorem{definition}{Definition}
\newtheorem{assumption}{Assumption}

%%%


\title{Title}


% The \author macro works with any number of authors. There are two commands
% used to separate the names and addresses of multiple authors: \And and \AND.
%
% Using \And between authors leaves it to LaTeX to determine where to break the
% lines. Using \AND forces a line break at that point. So, if LaTeX puts 3 of 4
% authors names on the first line, and the last on the second line, try using
% \AND instead of \And before the third author name.


\author{
  Name Surname\\
  MIPT\\
  Moscow, Russia\\
  \texttt{email@phystech.edu}\\
  \And
  Name Surname\\
  MIPT\\
  Moscow, Russia\\
  \texttt{email@phystech.edu}\\
}


\begin{document}


\maketitle

\begin{abstract}
    В данной работе решается задача детекции строк на фотографиях страниц рукописного текста на русском языке. Основная сложность заключается в обработке изображений с высоким уровнем шума, неравномерным, слабым освещением, геометрическими искажениями, а также наличием посторонних предметов на фотографии. Дополнительной проблемой является расположение строк, отличное от последовательности горизонтальных линий. В результате исследования планируется разработать новый алгоритм, который качественно превосходить существующие. Разработанные методы могут найти применение, например, в задачах распознавания древних исторических текстов.
\end{abstract}

\textbf{Keywords:} TODO

\section{Introduction}\label{sec:intro}

Детекция строк зачастую является необходимым этапом при распознавании рукописных текстов — так как модели распознавания работают на уровне строк. Она заключается в выделении на изображении отдельных элементов — текстовых строк. Качество работы всей системы распознавания рукописного текста напрямую зависит от результатов этого этапа: ошибки в выделении строк неизбежно приводят к снижению точности распознавания. Таким образом, разработка эффективного и надёжного метода детекции строк представляет собой актуальную задачу.

На пути к решению задачи возникает ряд проблем. Основными из них являются: непостоянство наклона строк текста в изображении, взаимное наложение строк друг на друга, неудовлетворительные условия освещения.

\subsection{Постановка задачи}
\subsubsection{Дано}
 Набор изображений страниц с рукописным текстом на русском языке: $\mathcal{I} = \{I_k\}_{k=1}^K$, где каждое $I_k \in \mathbb{R}_+^{H_k \times W_k \times C}$, $H_k$ — высота, $W_k$ — ширина, $C$ — количество каналов.
Для каждого изображения $I_k$ задано множество: $\mathcal{R}_k = \{R_{k,1}, \dots, R_{k, n_k}\}$, 
    где каждое $R_{k,j}$ — множество точек, которые задают минимальный ограничивающий прямоугольник $j$-й строки текста, а $n_k$ — количество строк на изображении. Под строкой текста понимается последовательность рукописных символов и слов, расположенных вдоль одной локально гладкой
    траектории письма и отделённая от соседних строк межстрочным промежутком.
Датасет: $\mathcal{D} = \{(I_k, \mathcal{R}_k)\}_{k=1}^K$.

\subsubsection{Найти}

Построить отображение
    \[
        f \colon \mathcal{I} \to 2^{B},
    \]
    где
    \[
    B =
    \left\{
    \{p_i\}_{i=1}^{4}
    \;\middle|\;
    p_i=(x_i,y_i)\in \mathbb{Z}_{\ge 0}^{2},
    \text{ точки } p_1,p_2,p_3,p_4
    \text{ являются вершинами прямоугольника}
    \right\}.
    \]
    
    По входному изображению $I \in \mathcal{I}$ отображение $f$ предсказывает набор минимально ограничивающих прямоугольников:
    \[
        f(I) = \hat{\mathcal{R}} \subseteq 2^{B}.
    \]

Необходимо, чтобы для любого предсказанного прямоугольника
    \[
        R = \{p_i\}_{i=1}^{4} \in \hat{\mathcal{R}},
        \qquad p_i=(x_i,y_i),
    \]
    выполнялось условие принадлежности всех его вершин области изображения:
    \[
        p_i \in \Omega_I,
        \qquad
        \Omega_I = \{0,\dots,W\}\times\{0,\dots,H\},
        \qquad i=1,\dots,4.
    \]
    Здесь $H$ и $W$ — высота и ширина изображения $I$.

\subsubsection{Критерий}
Для каждой пары предсказанного и истинного полигонов считается Intersection over Union:
\[
\operatorname{IoU}(\hat{R}_i, R_j)
=
\frac{
    \left|\hat{R}_i \cap R_j\right|
}{
    \left|\hat{R}_i \cup R_j\right|
}.
\]
Сопоставление выполняется жадно. Для каждого предсказанного полигона $\hat{R}_i$
выбирается ещё не занятый ground truth полигон $R_j$ с максимальным значением IoU:
\[
j^*
=
\arg\max_{j \in \mathcal{R}_{\text{free}}}
\operatorname{IoU}(\hat{R}_i, R_j).
\]

Если
\[
\operatorname{IoU}(\hat{R}_i, R_{j^*}) \ge 0.5,
\]
то такая пара считается true positive (TP).

Если подходящего ground truth полигона не найдено, предсказание считается false positive (FP).

Все ground truth полигоны, которые не были сопоставлены ни с одним предсказанием,
считаются false negative (FN).

Итоговые метрики считаются следующим образом:
\[
\operatorname{precision}
=
\frac{TP}{TP + FP},
\]
\[
\operatorname{recall}
=
\frac{TP}{TP + FN},
\]
\[
\operatorname{hmean}
=
\frac{
    2 \cdot \operatorname{precision} \cdot \operatorname{recall}
}{
    \operatorname{precision} + \operatorname{recall}
}.
\]

\subsubsection{Литература}

В литературе представлен ряд подходов для решения данной задачи. Классические методы машинного обучения, такие как метод горизонтальных проекций \cite{manmatha-2005}, преобразование Хафа \cite{likforman-1995}, алгоритмы размазывания \cite{shi-2004}, чаще всего решают задачу, не учитывая все вышеперечисленные проблемы. Только их комбинирование с другими методами может привести к приемлемому результату \cite{Louloudis-2008}.

Также есть нейросетевые подходы, которые всё чаще \cite{louloudis-2025} используются для решения данной задачи. \cite{Liao2022} Используются трансформеры для нахождения матриц вероятности и порогов для последующей бинаризации и выделения текста. Нейросетевые методы, обучение без учителя такие как UTLS \cite{barakat-2020-utls} очень чувствителен к сильной неоднородности структуры документа.


В современной литературе почти никогда не учитывается тип языка при распознавании строк, и работа ведётся с хорошо освещёнными отсканированными изображениями.

\section{Related Work}\label{sec:rw}

Метод UTLS \cite{barakat-2020-utls} чувствителен к сильной неоднородности структуры документа: большим различиям в высоте строк, межстрочных расстояниях, наклонах и форме строк. Это связано с тем, что размер анализируемого патча выбирается относительно средней высоты символов или строк, поэтому при сильных перепадах масштаба один и тот же размер окна может быть подходящим для одних строк и неподходящим для других.

Метод Kurar Barakat \cite{barakat-2020-learning-free} демонстрирует высокое качество на исторических рукописных документах: на DIVA—HisDB авторы сообщают достижение Line IU до 100\%, а на cBAD получают результаты, близкие к лучшим learning—free подходам, хотя уступающие лучшему learning—based методу с F—measure 97.70. При этом метод является вычислительно затратным: полный прогон по cBAD занимает около одного дня. Основные ограничения подхода связаны с высокой зависимостью от оценки высоты символов, ухудшением качества при сильном наклоне строк, появлением ложных blob lines и неточным разделением касающихся символов соседних строк.


\section{Preliminaries}\label{sec:prelim}


\section{Experiments}\label{sec:exp}

\subsubsection{Преобразование Хафа}

Наша команда пробовала метод, основанный на преобразовании Хафа \cite{Louloudis-2008}, данный подход имеет ряд недостатков.

Во-первых он очень зависит от качества бинаризации. Если, к примеру, какой-то метод бинаризации после себя оставит ряд мелких чёрных пикселей, это усложнит задачу подсчёта средней ширины или высоты компонент (смотри рис.~\ref{fig:small_components_bad_binarization}).

Во-вторых иногда несколько слов в идущих подряд строках могут образовывать одну компоненту связности, метод в статье не умеет разделять такие компоненты, а также в голосовании Хафа все они не будут участвовать что создаст сильную неоднородность количества голосующих точек между строками (смотри рис.~\ref{fig:connected_component_several_lines}).

В-третих необходимо подбирать гиперпараметры: порог голосов для основной линии порог для дополнительных линий радиус окрестности для объединения голосов порог отклонения угла для отбрасывания строк и многие другие.

В-четвёрых метод поиска, предложенный авторами, дополнительных строк значительно понижает precision и лишь немного увеличивает recall.

\begin{figure}[h!]
\centering

\begin{subfigure}{0.48\textwidth}
    \centering
    \includegraphics[width=\textwidth]{hough_method/Пример_множества_маленьких_компанент_при_плохой_бинаризации.png}
    \caption{Пример множества маленьких компонент при плохой бинаризации.}
    \label{fig:small_components_bad_binarization}
\end{subfigure}
\hfill
\begin{subfigure}{0.48\textwidth}
    \centering
    \includegraphics[width=\textwidth]{hough_method/Пример_связной_компаненты_слов_из_нескольких_подряд_идущих_строк.png}
    \caption{Пример связной компоненты слов из нескольких подряд идущих строк.}
    \label{fig:connected_component_several_lines}
\end{subfigure}

\end{figure}

Мы предлагаем следующие решения данных проблем:

1) Вместо выбора самой частой высоты брать среднее используя робастность, будем отсекать первый бин считая его шумом бинаризации (смотри рис.~\ref{fig:first_bin_cut_histogram}).

2) Мы предлогаем для начала разделить \text{subset\_2}. Это даст возможность всем словам участвовать в голосовании. Для этого мы смотрим на все связные компаненты из \text{subset\_2} и с помощью метода, описанного в статье \cite{das-2023-seam-carving} предварительно предварительно максимизировав дисперсию HPP с помощью поворота разделяем слова. Добавляем получившиеся компаненты в \text{subset\_1} и только после этого проводим голосование Хафа (смотри рис.~\ref{fig:subset2_split_example}).

3) Наша команда не смогла найти полное решение данной проблемы. Однако ряд гиперпараметров возможно заменить на менее чувствительные а именно диапазон углов в матрице голосования можно задавать так: сначала оценивать глобальный угол наклона строк с помощью максимизации дисперсии hpp потом с помощью гиперпараметра задавать отклонения относительно этого угла, это сделает метод гораздо устойчивее. Также был установлен один гиперпараметр для слияния линий и добавления новых на который домножается среднее значение расстояния между предсказанными линиями в проекции на вертикаль посреди изображения. Была введена проверка для новых линий что их угол лежит в диапазоне углов в матрице голосования

4) Наша команда решила убрать данную стадию, учитывая тот факт, что датасет состоит из сочинений и строки на фотографиях имеют примерно одинаковую длину.

\begin{figure}[h!]
\centering

\begin{subfigure}{0.85\textwidth}
    \centering
    \includegraphics[width=\textwidth]{hough_method/Пример_разделения_sub_2.jpg}
    \caption{Пример разделения \texttt{subset\_2}.}
    \label{fig:subset2_split_example}
\end{subfigure}

\vspace{0.5cm}

\begin{subfigure}{0.85\textwidth}
    \centering
    \includegraphics[width=\textwidth]{hough_method/Пример_гистограмма_отсечение_первого_бина.jpg}
    \caption{Пример гистограммы с отсечением первого бина.}
    \label{fig:first_bin_cut_histogram}
\end{subfigure}

\end{figure}

\subsubsection{Сравнение метода из статьи и нашего метода}
Наша команда провела эксперимент на датасете \cite{potyashin2023hwr200}. Были взяты 1569 изображений в качестве тестовой выборки из них было выделено 20, с помощью Optuna был произведён перебор гиперпараметров для метода (40 шагов, оптимизатор: optuna.samplers.TPESampler). На всех 1569 изображениях мы сравнили метод авторов и наш с внесёнными вышеописанными изменениями, в качестве предобработки в обоих случаях использовались модели \texttt{U—Net} и \texttt{Yolo} обученные на бинаризацию и сегментацию страниц текста соответственно.

\begin{table}[h!]
\centering
\caption{Сравнение метрик качества детекции строк}
\begin{tabular}{lccc}
\hline
Метод & Precision & Recall & Hmean \\
\hline
Paper method & 0.4519 & 0.5472 & 0.4950 \\
Our method   & 0.7426 & 0.7955 & 0.7681 \\
\hline
\end{tabular}
\end{table}

\subsubsection{Вывод}

В целом метод имеет большое количество гиперпараметров, что делает его применение сложным.

\subsubsection{Горизонтальный профиль проекции}

Мы также пробовали метод, основанный на горизонтальном профиле проекции \cite{das-2023-seam-carving}, данный подход имеет ряд недостатков:

1) Метод вообще не учитывает наклон строк даже глобальный (смотри рис.~\ref{fig:hpp_bad_without_global_skew_image} рис.~\ref{fig:hpp_bad_without_global_skew_energy}).

2) Анализ профиля проекции производится крайне поверхностно, задан один гиперпараметр порогового значения. Факт одинаковости строк по длине вовсе не гарантирует, что пики в профиле проекции будут примерно одинаковыми. В какой-то строке могут быть большие промежутки между словами, более размашистый почерк.

\begin{figure}[h!]
\centering

\begin{subfigure}{0.55\textwidth}
    \centering
    \includegraphics[width=\textwidth]{hpp_method/Пример_hpp_работает_плохо_без_исправления_глобального_наклона_изображение.jpg}
    \caption{Пример, где HPP работает плохо без исправления глобального наклона: исходное изображение.}
    \label{fig:hpp_bad_without_global_skew_image}
\end{subfigure}

\vspace{0.5cm}

\begin{subfigure}{0.55\textwidth}
    \centering
    \includegraphics[width=\textwidth]{hpp_method/Пример_hpp_работает_плохо_без_исправления_глобального_наклона_энергетическая_матрица.jpg}
    \caption{Пример, где HPP работает плохо без исправления глобального наклона: энергетическая матрица.}
    \label{fig:hpp_bad_without_global_skew_energy}
\end{subfigure}

\vspace{0.5cm}

\begin{subfigure}{0.55\textwidth}
    \centering
    \includegraphics[width=\textwidth]{hpp_method/Пример_работы_после_выранвнеивания_и_корекции_глобального_наклона_энергетическая_матрица.jpg}
    \caption{Пример работы после выравнивания и коррекции глобального наклона: энергетическая матрица.}
    \label{fig:hpp_after_alignment_global_skew_energy}
\end{subfigure}

\end{figure}

Наше улучшение методов на основе горизонтального профиля проекции:

1) Учёт глобального и локального наклона строк (смотри рис.~\ref{fig:hpp_after_alignment_global_skew_energy})

Перед вычислением локальных углов $\alpha(x, y)$ наша команда сначала исправляет \textbf{глобальный наклон страницы}:

1. Бинаризуем изображение $B(x, y)$:
\[
B(x, y) =
\begin{cases}
0, & \text{фон} \\
1, & \text{текст}
\end{cases}
\]  

2. Оцениваем глобальный угол $\theta_\text{global}$ методом горизонтальных проекций:
\[
\theta_\text{global} = \arg\max_{\theta \in [-30^\circ, 30^\circ]} \text{Var} \Big( \sum_y B_\theta(x, y) \Big)
\]
где $B_\theta$ — бинаризованное изображение, повернутое на угол $\theta$.

3. Применяем вращение изображения на угол $-\theta_\text{global}$:
\[
B_\text{corrected} = \text{rotate}(B, -\theta_\text{global})
\]

Далее делаем локальное выпрямление:

1. Делим изображение на сетку $G_x \times G_y$.  

2. Для каждой ячейки вычисляем локальный угол наклона $\alpha(x, y)$:
\[
\alpha(x, y) = \text{local\_angle}\Big( B_\text{cell}(x, y) \Big)
\]
где $\text{local\_angle}$ также оценивается через горизонтальный профиль внутри ячейки (аналогично глобальному методу).

3. Далее строим наклон $s(x, y) = \tan \alpha(x, y)$ и интегрируем вдоль горизонтали для корректного смещения пикселей:
\[
dy(x, y) = \sum_{t=0}^{x-1} s(t, y), \quad dy(x, y) \gets dy(x, y) - \frac{1}{W} \sum_{x=0}^{W-1} dy(x, y)
\]

4. Наконец, строим преобразование для выпрямления текста:
\[
T(x, y) = (x, y + dy(x, y)), \quad
B_\text{warp}(x, y) = B_\text{corrected}(T(x, y))
\]

2) Морфологический анализ горизонтального профиля (HPP)

После выпрямления профиля строки $H = [h_1, h_2, \dots, h_n]$ выполняется нормализация HPP:
\[
\tilde{H} = \frac{H - \min(H)}{\max(H) - \min(H)}
\]

Далее применяются операции \textbf{closing} и \textbf{opening} с ядром размера $k$:
\[
H_{\text{smoothed}} = \text{open}(\text{close}(\tilde{H}, k), k)
\]

\[
\text{open}(H, k)[y] = \max_{t \in S} \min_{s \in S, \, y+s+t \in [0,H—1]} H[y+s+t]
\]

\[
\text{close}(H,k)[y] = \min_{t \in S} \max_{s \in S, \, y+s+t \in [0,H—1]} H[y+s+t]
\]

\[
S = \{-\lfloor k/2 \rfloor, \dots, 0, \dots, \lfloor k/2 \rfloor\}
\]


И обратно масштабируется в исходный диапазон:
\[
H_{\text{final}} = H_{\text{smoothed}} \cdot (\max(H)-\min(H)) + \min(H)
\]

На сглаженном профиле $H_{\text{final}}$ определяется набор регионов так: мы идём по профилю если видим подъём на 5\% относительно текущего уровня и больше затем спуск аналогично на 5\% и более — это регион строки текста. Мы используем несколько размеров ядра $k_1, k_2$ (берём ядра, деля среднюю высоту посчитанную как описано выше на 2 и на 3) и выбираем вариант, дающий максимальное количество строк.

Был проведён аналогичный предыдущему эксперимент с предварительным подбором гиперпараметров. Результат следующий:

\begin{table}[h!]
\centering
\caption{Сравнение результатов двух методов HPP seam carving.}
\label{tab:hpp_seam_comparison}
\begin{tabular}{lccccc}
\hline
\textbf{Method} & \textbf{H-mean} & \textbf{Precision} & \textbf{Recall} \\
\hline
Paper method  & 0.7337 & 0.8492 & 0.6458 \\
Our method & 0.8419 & 0.8494 & 0.8345 \\
\hline
\end{tabular}
\end{table}

\subsubsection{Нейросетевые подходы}

Наша команда пробовала обучать нейросетевую модель описанную в статье \cite{Liao2022} на датасете \cite{potyashin2023hwr200} с предварительной разметкой моделей PaddleOCR-VL (1.5). Результаты по метрикам вышли следующие:

\begin{table}[h!]
\centering
\caption{Метрики качества DBNet++.}
\label{tab:dbnet_metrics}
\begin{tabular}{lccc}
\hline
\textbf{Method} & \textbf{H—mean} & \textbf{Precision} & \textbf{Recall} \\
\hline
DBNet++ & 0.9577 & 0.9805 & 0.9359 \\
\hline
\end{tabular}
\end{table}

\section{Discussion}\label{sec:disc}

TODO

\section{Conclusion}\label{sec:concl}

TODO


%%%%%%%%%%%%%%%%%%%%%%%%%%%%%%%%%%%%%%%%%%%%%%%%%%%%%%%%%%%%

\bibliographystyle{unsrtnat}
\bibliography{references}

%%%%%%%%%%%%%%%%%%%%%%%%%%%%%%%%%%%%%%%%%%%%%%%%%%%%%%%%%%%%

\newpage
\appendix
\section{Appendix / supplemental material}\label{app}

\subsection{Additional experiments / Proofs of Theorems}\label{app:exp}

TODO

\end{document}